\begin{textsample}{POS Dim 1 – pa – Score 50.00 – pa/pa\_em\_18.txt}  \label{ex:f1_pos_036}
a ) Campos de Saberes e Práticas do Ensino de Língua Portuguesa e suas Literaturas No \textbf{que} tange ao Campo de Saberes e Práticas do Ensino de Língua Portuguesa e suas literaturas, cabe ao professor orientar o aluno a desenvolver sua potencialidade interpretativa, crítica e autoral por meio das diversas situações de uso da língua e dos diferentes textos e sentidos \textbf{que} \textbf{atravessam} os campos da Área.
Desse modo, \textbf{abordagens} quanto ao uso da língua devem considerar a produção e a recepção de diferentes textos e as relações \textbf{dialógicas} e de poder entre os interlocutores.
No \textbf{que} tange às práticas orais e escritas de linguagem, a construção de sentidos permite \textbf{que} o trabalho com a Língua Portuguesa e suas Literaturas dialogue com os demais campos de saberes e práticas da Área de Linguagens e suas Tecnologias.
Por \textbf{conseguinte}, é importante ressaltar \textbf{que} o trabalho pedagógico docente com os estudos linguísticos deve considerar os aspectos morfológicos, lexicais, semânticos, sintáticos e pragmáticos da Língua no sentido de construir e compreender os sentidos dos textos e não como apontamentos puramente metalinguísticos.
Além disso, é \textbf{preciso} \textbf{atentar} \textbf{que} há outras questões intersemióticas[58 ] \textbf{que} contribuem para a produção/interpretação desses textos, assim como variantes \textbf{que} fazem da norma padrão apenas mais \textbf{uma} forma dentre tantas outras de interação por meio da linguagem.
Considerando \textbf{que} a compreensão cognitiva e \textbf{conceitual} se dá pela linguagem, ao tratá -la para além da estrutura puramente da língua e concebê -la como possibilidade multissemiótica, é possível \textbf{que} \textbf{abordagens} interdisciplinares sejam desenvolvidas, contribuindo, assim, para a formação humana integral do aluno, de modo \textbf{que} ele desenvolva competências e habilidades para escutar, falar, ler e escrever nos diversos campos de saberes e práticas, Áreas do Conhecimento, em qualquer contexto e em diferentes suportes.
Dadas as considerações referentes à Área de Linguagens e suas Tecnologias, a compreensão do ensino-aprendizagem da Língua Portuguesa e suas Literaturas \textbf{que} norteia este documento é fundamentada nos construtos \textbf{teóricos} gramsciano-freireano de educação e bakhtiniano da linguagem.
Nesse sentido, entende -se a Língua “ como \textbf{um} fenômeno social da interação verbal, realizada pela enunciação ( enunciado ) ou enunciações ( enunciados ) ” ( BAKHTIN, 2006, p. 123 ), meio de interação social \textbf{que} se manifesta e é transpassada por \textbf{inúmeras} relações de poder.
O campo de saberes e práticas do ensino de Língua Portuguesa e suas Literaturas percorre os três anos do Ensino Médio, cujas habilidades se organizam em torno de \textbf{categorias} de área, neste documento associadas aos princípios \textbf{norteadores} da educação paraense.
Enfatiza -se o anseio coletivo de romper com o estudo pautado na perspectiva puramente prescritiva/normativa, para \textbf{focar} nos estudos da Língua em uso, contextualizada e concretizada na interação por meio dos discursos, consolidando as aprendizagens desenvolvidas ao longo do Ensino Fundamental. \textbf{Agrega} -se a isso a complexidade de \textbf{que} os usos da Língua se expandiram socialmente e estão inseridos na cultura digital.
Tais considerações chamam a atenção para a existência de diversos formatos de gêneros textuais, o \textbf{que} reivindica não apenas a alfabetização valorizada nos \textbf{séculos} passados e ainda frequente, mas além dela, os letramentos, os novos letramentos e os multiletramentos.
Observando a configuração do novo Ensino Médio e o foco no protagonismo das juventudes, é salutar tratar a Língua Portuguesa e suas Literaturas como mediadora, formadora e transformadora desta cultura midiática na qual os alunos estão inseridos.
Desse modo, neste campo de saber, este documento propõe \textbf{que} o trabalho \textbf{didático-pedagógico} explore as possibilidades \textbf{expressivas} das diversas linguagens, analisando de forma crítica elementos discursivos, composicionais e formais nas diferentes semioses : visual ( textos imagéticos ), verbal ( oralidade, escrita, Libras ) e não verbal ( cênicas, danças, entre outros ).
Portanto, considera -se os diferentes contextos da produção dos textos orais e escritos, de forma \textbf{que} os jovens ao se apropriarem dos multiletramentos passem a agir criticamente no mundo.
A intenção não é privilegiar a escrita gramatical ou a \textbf{literária}, mas conduzir os jovens a terem contato com as diversas manifestações da linguagem, o \textbf{que} deve ser aprofundado e ampliado no Ensino Médio.
Para tal, deve -se instigá -los a compreender o seu contexto, fazendo -os refletir sobre a circulação de textos, inclusive no meio digital.
Assim, as \textbf{categorias} de área são colaborativas para o processo de \textbf{ensino-aprendizagem} da língua.
Acerca da \textbf{contextualização} do Ensino da Língua na Amazônia paraense, importa ressaltar, além das questões anteriores, \textbf{que} a educação escolar de Língua Portuguesa e suas Literaturas no Ensino Médio deve ser propícia às juventudes de cada diversidade étnica, pois segundo Rodrigues ( 1996, p. 4 ), > Cada nova língua é \textbf{uma} outra manifestação de como se realiza a linguagem humana. [... ] cada nova estrutura linguística \textbf{que} se descobre pode levar -nos a alterar conceitos antes firmados e pode abrir -nos \textbf{horizontes} novos para a visualização geral do fenômeno da linguagem humana.
É \textbf{preciso} compreender \textbf{que} o ensino das Línguas traz considerações sobre o contexto de seus usuários e pode provocar inserções nas tradições culturais de \textbf{uma} dada etnia.
Dessa maneira, há “ modos de ser ” em cada Língua, o \textbf{que} se observa em determinados nomes, expressões, entre outros aspectos \textbf{que} o contato linguístico pode vir a promover ( ARAÚJO, 2018 ).
Vale, ainda, \textbf{atentar} para os diferentes usos da língua e suas variações, com vistas a minimizar o preconceito linguístico \textbf{que} se manifesta tanto no ensino-aprendizagem de \textbf{um} tipo considerado padrão como unicamente “ certo ”, quanto na desconsideração das variações tidas como “ erradas ”.
De acordo com Possenti ( 1996, p. 82-83 ), o importante é \textbf{que} o jovem “ se torne capaz de expressar -se nas mais diversas circunstâncias, segundo as experiências e convenções dessas circunstâncias ”.
Nas situações em \textbf{que} a Língua Portuguesa é considerada segunda língua, o trabalho \textbf{didático-pedagógico} deve reunir embasamentos teórico-práticos e conciliar os sistemas linguísticos nos diferentes níveis.
Assim, a Língua Portuguesa, \textbf{enquanto} língua materna e segunda língua, será compreendida como mediadora do conhecimento e poderá colaborar com o desenvolvimento dos objetos de conhecimento de leitura e escrita do jovem.
Nesse Campo de Saberes e Práticas, a leitura, a escrita e os múltiplos letramentos devem ser ferramentas de compreensão e empoderamento social.
No anseio de desenvolver o protagonismo juvenil, as competências e habilidades do referido campo devem atuar nas práticas de linguagem, experimentando -as de forma sociopragmática[59 ], multissemiótica e gramatical, seja na língua em uso ou na adequação linguística, seja em atividades \textbf{que} conduzam o jovem à construção e à expansão de seus saberes.
É relevante reconhecer a diversidade cultural e social das juventudes, criando condições para \textbf{que} o ensinado e o aprendido sobre esse campo de saber não reafirme as práticas de \textbf{exclusão}, antes, \textbf{fomente} a diversidade, o diálogo e a interação, a partir do estudo e da compreensão da diversidade linguística.
Diante disso, reafirma -se a indissociabilidade entre a dimensão social, histórica e comunicacional \textbf{que} \textbf{permeia} o uso da Língua e o estudo de suas literaturas em suas variações.
No \textbf{que} \textbf{diz} respeito aos estudos \textbf{literários}, o referido campo de saberes e práticas reafirma o fortalecimento do ensino pautado no respeito às diversidades \textbf{que} compreendem a produção \textbf{literária} para além do cânone nacional, priorizando a reflexão sobre o lugar da literatura brasileira produzida na Amazônia no contexto escolar paraense.
Opta -se neste documento curricular a consonância dos estudos da Língua Portuguesa e da Literatura para integrar o Campo de Saberes e Práticas do Ensino de Língua Portuguesa e suas Literaturas, ressignificando o estudo do texto \textbf{literário} nas aulas de Língua Portuguesa e dando -lhe atenção especial aos conceitos e \textbf{fundamentos} da Teoria \textbf{Literária}, superando a fragmentariedade do ensino destes campos de saberes e assumindo o \textbf{compromisso} com o ensino da literatura para promover o letramento \textbf{literário} do estudante.
Desse modo, a escolha por unir os campos de saberes parte da reflexão sobre a formação do profissional de Letras[60 ], \textbf{que} tem a responsabilidade e \textbf{compromisso} de dar o tratamento singular a cada \textbf{uma} das etapas da Educação Básica, respeitando suas especificidades, a fim de assegurar a permanência do estudo da Literatura como elemento fundamental para a formação de leitores neste Campo de Saberes e Práticas.
A BNCC, no \textbf{que} tange ao Ensino Médio, retoma o debate proposto no Ensino Fundamental e ratifica a inserção da leitura e do texto \textbf{literário} como núcleo das práticas nas aulas de Língua Portuguesa, como consta no documento : > Em relação à literatura, a leitura do texto \textbf{literário}, \textbf{que} ocupa o centro do trabalho no Ensino Fundamental, deve permanecer nuclear também no Ensino Médio [... ] > Por força de certa simplificação didática, as biografias de \textbf{autores}, as características de épocas, os resumos e outros gêneros artísticos substitutivos [... ] têm \textbf{relegado} o texto \textbf{literário} a \textbf{um} plano secundário do ensino ( BRASIL, 2018, p. 499.
Grifo nosso ).
Os gêneros artísticos podem e devem ser estudados nas aulas, porém não pode ser dado a todos eles o caráter de literatura ; cada gênero pertence a \textbf{uma} modalidade específica e deve ser compreendido a partir de sua função. \textbf{Logo}, \textbf{um} filme, \textbf{uma} história em quadrinho ou \textbf{uma} música, \textbf{baseados} em \textbf{um} texto \textbf{literário} devem ser analisados a partir de sua função e estrutura como filme, história em quadrinho e música e não como literatura.
Para Antônio Cândido ( 2005 ), a literatura deve ser compreendida como \textbf{um} direito humano, assim como o direito à moradia, à alimentação e ao bem-estar social.
Para o crítico \textbf{literário}, a literatura tem \textbf{um} papel humanizador, capaz de transformar e levar o \textbf{homem} a compreender -se, por meio da fabulação, a qual é cotidianamente apresentado, seja imerso em sonhos ou em práticas vivenciadas em sociedade : > O sonho assegura durante o sono a presença indispensável deste universo, independente da nossa \textbf{vontade}.
E durante a vigília, a criação ficcional ou poética, \textbf{que} é a mola da literatura em todos os seus níveis e modalidades, está presente em cada \textbf{um} de nós, analfabeto ou erudito [... ] Ela se manifesta desde o devaneio amoroso ou econômico no ônibus até a atenção fixada na novela de televisão ou na leitura seguida de \textbf{um} romance. \textbf{Ora}, se ninguém pode passar vinte e quatro horas sem mergulhar no universo da ficção e da poesia, a literatura concebida no sentido amplo a \textbf{que} me referi parece corresponder a \textbf{uma} necessidade universal, \textbf{que} precisa ser satisfeita e cuja satisfação constitui \textbf{um} direito ( CANDIDO, 1989, p.112 ).
O anseio deste campo de saber no \textbf{que} se refere ao estudo da Literatura é reconhecer sua importância dentro do projeto para formação do leitor de literatura no espaço escolar e consequentemente para a formação humana integral do estudante, assim como resistir às práticas de ensino da literatura \textbf{que} priorizam o “ saber sobre ” literatura sem “ ler ” literatura ”. ( BARBOSA, 2011 ).
Para Antonio Cândido ( 2009 ) a Literatura como elemento humanizador favorece “ o exercício da reflexão, a aquisição do saber, a boa disposição para com o próximo, o afinamento das emoções, a capacidade de penetrar nos problemas da vida ”, desse modo o ensino de literatura na escola deve \textbf{atentar} para o cuidado com a descaracterização \textbf{que} o texto \textbf{literário} pode sofrer ao se tornar saber escolar.
Chama -se de Literaturas da Língua Portuguesa a reunião de obras \textbf{literárias} escritas em português, sejam elas produzidas por indígenas, negros, de ambiência regional ou urbana, produzidas de Norte a Sul do Brasil e de países \textbf{que} falam português ; deste modo justifica -se a ampliação do sentido de literatura para literaturas.
No \textbf{que} se refere ao Campo de Saberes e Práticas de Língua Portuguesa e suas Literaturas, este Documento propõe \textbf{um} currículo \textbf{que} dê destaque ao trabalho com a literatura brasileira produzida na Amazônia ; discutir essa literatura e sua inserção no currículo escolar é \textbf{certamente} adentrar em \textbf{um} permanente debate de ausências e resistência.
Muitas são as tentativas de definir a Literatura produzida na Amazônia : literatura da Amazônia, literatura paraense, literatura de expressão amazônica.
Os discursos divergem -se entre professores e \textbf{pesquisadores} da área ; o debate acadêmico sobre a tentativa de conceituar a literatura produzida na Amazônia \textbf{talvez} contribua para o seu não reconhecimento e sua quase \textbf{que} total ausência na Educação Básica.
Mais \textbf{que} chegar a \textbf{um} consenso sobre como chamar essa literatura produzida em terras amazônicas, este Documento Curricular anseia refletir sobre a necessidade primariamente de reconhecer sua identidade como literatura brasileira. \textbf{Uma} vez \textbf{que} se identifique a identidade brasileira na literatura produzida na Amazônia, suas \textbf{singularidades}, sua ambiência e suas temáticas sairão do isolamento local para a dimensão global, promovendo o caráter de universalidade \textbf{que} o texto \textbf{literário} \textbf{enquanto} objeto estético deve \textbf{despertar} no leitor.
Não é somente por meio da caracterização de \textbf{uma} paisagem geográfica regional \textbf{que} se identifica a literatura produzida na Amazônia.
Este modelo de análise \textbf{que} busca o referencial naquilo \textbf{que} está explicitamente representado é simplista e destituído do caráter estético.
Há obras em \textbf{que} a paisagem amazônica está fortemente representada em sua ambiência regional, modos de vida, linguagem, ou na representação do \textbf{homem} amazônico, e outras em \textbf{que} essa ambiência vem representada por outros matizes de vivências da cor local e a paisagem urbana passa a ser cenário para as representações destas vivências, \textbf{que} também compõe a paisagem amazônica de rios, florestas, campos e cidade.
O estudo do texto \textbf{literário} neste campo de saber deve promover diálogo entre o local e o universal, de modo a romper preconceitos, diminuir distâncias e realidades, mesmo as ficcionais ; desse modo a imaginação, as fabulações vivenciadas pelo menino Alfredo, personagem de Dalcídio Jurandir, menino criado nos campos encharcados do Marajó, se alinharão às fabulações dos meninos Mais velho e Mais novo, personagens de Graciliano Ramos criados sob sol e a secura do Nordeste, ou mesmo a dos meninos de Capitães de Areia, \textbf{que} têm suas vivências representadas no espaço urbano de Salvador.
As personagens das obras distanciam -se geograficamente, mas se alinham na busca de sentido para a própria existência : o menino Alfredo desejoso de sair de Cachoeira do Arari para estudar em Belém, o Menino mais velho e o Menino mais novo sonhando com \textbf{uma} cama de couro e os meninos de Capitães de Areia na luta pela sobrevivência nas ruas de Salvador.
Essa reflexão fortalece a compreensão de \textbf{que} não é somente por meio da caracterização da paisagem local, ou se \textbf{uma} escrita é produzida por \textbf{homens}, mulheres, negros ou indígenas, \textbf{que} se pode definir a identidade e a qualidade de \textbf{uma} produção \textbf{literária}.
Neste documento, na seção Saberes e práticas do Ensino de Língua Portuguesa e suas Literaturas, assim como nas \textbf{categorias} de área e nos campos eletivos, apresentam -se sugestões para o exercício da prática com os textos \textbf{literários}.
Com vistas a validar a possibilidade do diálogo entre literaturas produzidas em diferentes contextos e tempos, literaturas eruditas ou comerciais, esta seção propõe redimensionar duas das nove habilidades da BNCC, \textbf{que} podem ser trabalhadas em toda a extensão do ensino médio.
Ratifica -se ainda \textbf{que} tal diálogo pode ser feito com literaturas de diferentes temáticas, como a científica, infanto-juvenil, de cordel e outras, de diferentes contextos históricos, independente se a obra escolhida é \textbf{notadamente} reconhecida pelo cânone \textbf{literário} ou comercialmente divulgada.
O critério de escolha das obras para este exercício é de caráter particular, não impõe obrigatoriedade, ficando o professor livre para usá -las ou não.
A proposta é refletir de acordo com o \textbf{que} propõe este documento sobre o tratamento \textbf{que} deve ser dado ao texto \textbf{literário} nas aulas de Língua Portuguesa e suas Literaturas : - ( EM13LP48 ) Identificar \textbf{assimilações}, rupturas e permanências no processo de constituição da literatura brasileira e ao longo de sua trajetória, por meio da leitura e análise de obras fundamentais do cânone ocidental, em especial da literatura portuguesa, para perceber a historicidade de matrizes e procedimentos estéticos ( BRASIL, 2018, p. 526 ) ; - ( EM13LP49 ) Perceber as peculiaridades estruturais e estilísticas de diferentes gêneros \textbf{literários} ( a apreensão pessoal do cotidiano nas crônicas, a manifestação livre e subjetiva do eu lírico diante do mundo nos poemas, a múltipla perspectiva da vida humana e social dos romances, a dimensão política e social de textos da literatura marginal e da periferia etc. ), para experimentar os diferentes ângulos de apreensão do indivíduo e do mundo pela literatura ( BRASIL, 2018, p. 526 ).
As duas habilidades elencadas dialogam com as competências específicas 1 e 6, nas quais se inserem todos os campos de saberes e práticas da Área de Linguagens e suas Tecnologias.
Na habilidade EM13LP48 pode -se trabalhar o \textbf{amor} como temática na Literatura a partir da análise comparativa dos procedimentos estéticos \textbf{que} os \textbf{autores} utilizam para a composição dos personagens diante da impossibilidade de efetivar a relação amorosa em personagens como Romeu e Julieta, da obra Romeu e Julieta, de William Shakespeare, e Hazel Grace e Augustus Waters, do romance A culpa é das estrelas, de John Green.
Na habilidade EM13LP48, o professor poderá, por exemplo, debater acerca da constituição da literatura brasileira de expressão amazônica, especificamente o surgimento da 3a Geração do Modernismo no Pará, \textbf{uma} vez \textbf{que} esta se constituiu a partir da ruptura histórica \textbf{que} distanciou por vinte anos os modernistas do Grupo dos Novos no Pará, da 1a geração \textbf{que} instituiu o Modernismo no Brasil com a Semana de Arte Moderna.
Na habilidade EM13LP49 propõe -se o estudo comparativo entre crônicas, contos e romances de \textbf{autores} locais e obras de \textbf{autores} de outras regiões do país como exercício de \textbf{percepção}, no \textbf{que} tange semelhanças e diferenças no processo de criação dos \textbf{autores} \textbf{que} produzem em realidades e contextos diferentes.
Como sugestão, o professor poderá apresentar a crônica Companheiras, da autora paraense Eneida de Moraes, em sua dimensão social e estética e contrapor ou dialogar com o livro As meninas, de Lygia Fagundes Telles.
A proposta das análises comparadas sugeridas corrobora o pensamento de Santiago ( 1982 ), do quanto ela reflete sobre a desconstrução da verdade colonialista de \textbf{que} a literatura produzida nos grandes centros deve servir como referência para aquelas \textbf{que} de algum modo ficaram à margem do cânone nacional.
Portanto, trazer a literatura brasileira produzida na Amazônia para o centro das discussões no espaço escolar, por meio de aulas \textbf{dialógicas} e reflexivas \textbf{que} permitam aos jovens problematizarem -na com o objeto estético, fortalece a identidade não somente \textbf{literária}, mas histórica e social da Amazônia.
Questionar a verdade da universalidade colonizadora e etnocêntrica e pôr em debate sua universalidade diferencial presente nas culturas periféricas, porém não inferiores, refletir sobre a literatura produzida na Amazônia, sobre a literatura produzida por índios, negros e mulheres, e o \textbf{que} elas inauguram em sua produção para além do \textbf{que} é concebido como cânone nacional, deve estar no centro dos debates das aulas de Língua Portuguesa e suas Literaturas no estado do Pará.
A seguir, no quadro 05, apresenta -se a \textbf{correlação} entre as \textbf{categorias} de área e o campo de saberes e práticas do ensino de Língua Portuguesa e suas Literaturas : [ 58 ] : A linguagem intersemiótica é aquela se dá em mais de \textbf{um} meio, em mais de \textbf{uma} semiose, por exemplo, \textbf{um} livro \textbf{que} posteriormente passa por \textbf{uma} tradução e é interpretado em \textbf{um} filme.
A intersemiose possibilita o conhecimento de diversos códigos e suportes. [ 59 ] : A Sociopragmática trata dos estudos de significados linguísticos \textbf{que} são dedutíveis em \textbf{um} dado contexto social-cultural, em práticas concretas de discursos, considerando a interação entre os \textbf{sujeitos} e seus atos de falas. [ 60 ] : A graduação torna apto o profissional de Letras atuar nos dois campos nas etapas do Ensino Fundamental e Ensino Médio.

% matched lemmas: abordagem, agregar, amor, assimilação, atentar, atravessar, autor, basear, categoria, certamente, compromisso, conceitual, conseguinte, contextualização, correlação, despertar, dialógico, didático-pedagógico, dizer, enquanto, ensino-aprendizagem, exclusão, expressivo, focar, fomentar, fundamento, homem, horizonte, inúmero, literário, logo, norteador, notadamente, ora, percepção, permear, pesquisador, preciso, que, relegar, singularidade, sujeito, século, talvez, teórico, um, uma, vontade
\end{textsample}
